\documentclass[12pt,a4paper]{article}
\usepackage{cmap} % Makes the PDF copiable. See http://tex.stackexchange.com/a/64198/25761
\usepackage[T1]{fontenc}
\usepackage[brazil]{babel}
\usepackage[utf8]{inputenc}
\usepackage{amsmath}
\usepackage{amsfonts}
\usepackage{amssymb}
\usepackage{amsthm}
\usepackage{textcomp} % \degree
\usepackage{gensymb} % \degree
\usepackage[usenames,svgnames,dvipsnames]{xcolor}
\usepackage{hyperref}
\usepackage{multicol}
\usepackage{graphicx}
\usepackage[margin=2cm]{geometry}

\hypersetup{
    colorlinks = true,
    allcolors = {blue}
}

% TODO: Consider using exsheets
% http://linorg.usp.br/CTAN/macros/latex/contrib/exsheets/exsheets_en.pdf
%
% http://ctan.org/tex-archive/macros/latex/contrib/exercise/
% Options: answerdelayed,lastexercise,noanswer
\usepackage[answerdelayed,lastexercise]{exercise}

\addto\captionsbrazil{%
\def\listexercisename{Lista de exerc\'icios}%
\def\ExerciseName{Exerc\'icio}%
\def\AnswerName{Solu\c{c}\~ao do exerc\'icio}%
\def\ExerciseListName{Ex.}%
\def\AnswerListName{Solu\c{c}\~ao}%
\def\ExePartName{Parte}%
\def\ArticleOf{de\ }%
}

\renewcommand{\ExerciseHeaderTitle}{(\ExerciseTitle)\ }
\renewcommand{\ExerciseListHeader}{%\ExerciseHeaderDifficulty%
\textbf{%\ExerciseListName\
\ExerciseHeaderNB.\ %
%\ --- \
\ExerciseHeaderTitle}%
%\ExerciseHeaderOrigin
\ignorespaces}
\renewcommand{\AnswerListHeader}{\textbf{\ExerciseHeaderNB.\ (\AnswerListName)\ }}

\newcommand{\fixme}{{\color{red}(...)}}

\newcommand*\sen{\operatorname{sen}}
\newcommand*\tg{\operatorname{tg}}
\newcommand*\R{\mathbb{R}}

\renewcommand{\theenumi}{\alph{enumi}}
\renewcommand\labelenumi{(\theenumi) }

\newcommand*\tipo{Prova III}
\newcommand*\turma{TADS121-01C}
\newcommand*\disciplina{CDI0001}
\newcommand*\eu{Helder G. G. de Lima}
\newcommand*\data{12/05/2017}

\author{\eu}
\title{\tipo - \disciplina}
\date{\data}

\begin{document}
\thispagestyle{empty}
\newgeometry{margin=2cm,bottom=0.5cm}
\begin{center}
\includegraphics[width=9.0cm]{marca} \\
\textbf{\tipo\ (\disciplina / \turma)} \\
Prof. \eu\footnote{
Este é um material de acesso livre distribuído sob os termos da licença \href{https://creativecommons.org/licenses/by-sa/4.0/deed.pt_BR}{Creative Commons Atribuição-CompartilhaIgual 4.0 Internacional}}
\end{center}

\noindent Nome do(a) aluno(a): \underline{\hspace{9,7cm}} Data: \underline{\data}

\begin{center}\fbox{
\begin{minipage}{14cm}

{\footnotesize
\begin{itemize}
\renewcommand{\theenumi}{\Roman{enumi}}
\item Identifique-se em todas as folhas.
\item Mantenha o celular e os demais equipamentos eletrônicos desligados durante a prova.
\item Justifique cada resposta com cálculos ou argumentos baseados na teoria estudada.
\item Resolva apenas os itens de que precisar para somar 10,0 pontos.
\end{itemize}
}

\end{minipage}
}
\end{center}

\begin{ExerciseList}
\Exercise[title={2,0}] Obtenha a derivada das seguintes funções (os cálculos apresentados devem deixar claro como encontrou cada resposta):
\begin{multicols}{3}
\begin{enumerate}
\item $f(x) = \sen(x) e^{(1-x^2)}$
\item $g(x) = \dfrac{(2x - 8)^{5}}{(x-3)^{6}}$
\item $h(x) = \ln( \cosh(3x-9) )$, sendo $\cosh(x) = \frac{e^x+e^{-x}}{2}$
\end{enumerate}
\end{multicols}
\Answer \fixme
\Exercise[title={2,0}] Seja $f(x) = \dfrac{ \tg(\pi x) }{ 2\pi }$. Calcule $f^\prime(0)$ usando a definição de derivada (isto é, um limite).
\Answer \fixme
\Exercise[title={2,0}] Determine o(s) valor(es) de $a \in \R$ sabendo que a função a seguir é diferenciável:
\[
p(x) =
\begin{cases}
\dfrac{1}{ax + 1/2}, \text{ se } x \geq 0,\\
\sqrt{4-x}, \text{ se } x \leq 0.
\end{cases}
\]
\Answer \fixme
\Exercise[title={2,0}] Verifique se a função $y = e^{-3x} + 2x e^{-3x}$ satisfaz a equação $y^{\prime\prime} + 6 y^\prime +9y = 0$.

\Answer \fixme
\Exercise[title={2,0}] Encontre a equação da(s) reta(s) paralela(s) a $r: y=6x+1$ que tangenciam o gráfico de $f(x) = 2x^3 - 6x^2 + 6x$.
\Answer \fixme
\Exercise[title={2,0}] Dê um exemplo de uma função $f: \R \to \R$ que seja contínua em algum ponto $a \in \R$ mas que não seja diferenciável em $a$. Apresente todos cálculos necessários para mostrar que $f$ é, de fato, contínua no ponto escolhido, e que não é derivável em tal ponto.
\Answer \fixme
\end{ExerciseList}

\begin{center}
BOA PROVA!
\end{center}

%\newpage
%\restoregeometry
%\section*{Respostas}
%\shipoutAnswer
\end{document}
